\section{Introduction}

A generated explanation can become misleading when a sensor window is corrected,
a feature fails a quality check, or a delayed observation replaces a statement
of absence. Deleting the answer discards useful information. Keeping it may
preserve a wrong number. A temporal evidence graph records what the answer
asserted, which observations and prerequisites it cited, and what was known when.

A comparison may cite both a parent claim and the two operands directly.
Revising either operand then reaches the comparison in one step.
If an answer relies only on an intermediate claim, however, a direct checker
may update that claim and leave the answer unchanged. A reached claim may also
remain supported because another valid route survives. Understanding these
cases requires the meaning of the edges and support rules, not visual complexity.

We study a typed temporal property graph whose schema is specified by the
researcher and whose instances record provenance and generated propositions.
First, we define how the graph is constructed
and how a query selects a semantic core linked to the stored objects.
A shared network view lets a reviewer inspect original citations, version changes
and groups of supporting inputs in the paper and an evidence-review dashboard.
Second, we connect an explicit support model to a locality result and a measurable
condition under which direct and transitive maintenance agree. Third, we test
that connection using actual Neo4j exports, independent proposition scores and
a controlled structural extension of three completed benchmarks. The graph
records commitments made by the answer, not the model's internal neural reasoning.

The waveform study provides a useful contrast. Only \SemParentClaims{}
of \SemDisplayedClaims{} claims displayed by transitive graph maintenance (M1)
declare parents, and direct checking
matches transitive maintenance. Every eligible waveform revision batch has zero
transitive exposure. A separate scoring protocol distinguishes accurate background
statements from violations of the original target-only output contract. It reveals
how checking trades coverage of required facts for grounded content. We then vary
dependency depth and alternative support routes in exact numerical programs and
common GPU-generated candidates over the same recordings. Together, these results
show when propagation matters and when direct maintenance suffices.

\subsection{Relation to Prior Work}

Truth-maintenance systems revise justified beliefs when their premises
change~\cite{doyle1979}. Incremental view maintenance updates database results
when their inputs change~\cite{gupta1993}. Provenance semirings distinguish
jointly required inputs from alternative derivations~\cite{green2007}.
We connect these established ideas to the dependencies a bounded generator
actually expresses and to independently measured maintenance outcomes.

PROV-DM distinguishes entities, activities, derivation, revision and
attribution~\cite{prov2013}. It informs our separation of recording provenance,
generated assertions and review events. We implement a property graph with
Python admission checks and database identity constraints. GraphRAG builds entity
graphs and community summaries for corpus-level question answering~\cite{edge2025}.
Graphiti, described in the Zep architecture, maintains temporal information for
agent memory~\cite{rasmussen2025}. Our unit of evaluation is a revisable numerical
proposition and its admissible supporting inputs. We measure whether evidence
grounds that proposition. This differs from asking whether a model's account of
its own computation faithfully describes its decision process~\cite{turpin2023}.
